\begin{lemma}[Blocking the action tensors]
    \label{lem:mpug_action_blocking}
    \lean{MPOTensor.blockTensor_actTensor,
        MPOTensor.GroupFamily.CarriesMPV.blockTensor,
        MPOTensor.GroupFamily.BlockActionData.block}
    \leanok
    \uses{def:mpug_permuted_block_l_symbol, def:mpug_action_tensor}
    For an operator tensor $T$ and a state tensor $A$, physical blocking
    satisfies $(T\cdot A)^{[L]}=T^{[L]}\cdot A^{[L]}$. A reduction
    $(V,W)$ from $T\cdot A$ onto $C$ remains a reduction from
    $T^{[L]}\cdot A^{[L]}$ onto $C^{[L]}$, with the same matrices.
    For $L>0$, the periodic action law
    $O_N(T)\ket{V^{(N)}(A)}=\ket{V^{(N)}(C)}$ at all positive lengths
    is also preserved by blocking.
\end{lemma}

\begin{proof}\leanok
    \uses{def:mpug_action_tensor, thm:mps_reduction_exterior_buffer}
    Decode a blocked input index as a word of length $L$. Reindexing
    the contracted input sum gives
    \begin{align}
        (T\cdot A)^{[L],I}
         & =\sum_J T^{[L],IJ}\otimes A^{[L],J}.
        \label{eq:mpug_action_blocking}
    \end{align}
    A word of blocked letters is the corresponding flattened word of
    original letters. The reduction equations and the periodic action law
    therefore apply at the flattened lengths.
\end{proof}

\begin{definition}[A common action buffer]
    \label{def:mpug_action_buffer}
    \lean{MPOTensor.GroupFamily.BlockActionData.nilpotencyLength,
        MPOTensor.GroupFamily.BlockActionData.HasExteriorBufferLength,
        MPOTensor.GroupFamily.BlockActionData.commonBufferLength,
        MPOTensor.GroupFamily.BlockActionData.commonBufferLength_pos,
        MPOTensor.GroupFamily.BlockActionData.nilpotencyLength_le_commonBufferLength_add_one,
        MPOTensor.GroupFamily.BlockActionData.commonBufferLength_of_isEmpty}
    \leanok
    \uses{def:mpug_permuted_block_l_symbol, thm:mps_reduction_exterior_buffer}
    Fix action reductions $(V_{g,x},W_{g,x})$ for finite $G$ and
    $\mathsf X$, and put $B_{g,x}=O_g\cdot A_x$. Write $N_0(g,x)$ for
    their residual nilpotency lengths. Their common buffer length is
    $L=\max\{1,\max_{g,x}(N_0(g,x)-1)\}$, where the inner maximum is
    zero for the empty family. Thus $L>0$ and $N_0(g,x)\leq L+1$;
    when $\mathsf X$ is empty, $L=1$.
\end{definition}

\begin{lemma}[Common blocking of action reductions]
    \label{lem:mpug_action_common_blocking}
    \lean{MPOTensor.GroupFamily.BlockActionData.interior,
        MPOTensor.GroupFamily.BlockActionData.exterior,
        MPOTensor.GroupFamily.BlockActionData.block_hasExteriorBufferLength_one,
        MPOTensor.GroupFamily.BlockActionData.block_commonBufferLength_hasExteriorBufferLength_one,
        MPOTensor.GroupFamily.exists_block_blockActionData_one,
        MPOTensor.GroupFamily.exists_block_blockActionData_one_of_isPretransitive}
    \leanok
    \uses{def:mpug_action_buffer, lem:mpug_action_blocking,
        thm:mpug_permuted_block_l_symbol_spec}
    Suppose $G$ and $\mathsf X$ are finite, the tensors $A_x$ are normal
    of positive, possibly different, bond dimensions, and the periodic
    action law holds at every positive length. There is one positive
    length $L$ for which all the blocked action tensors
    $B_{g,x}=O_g^{[L]}\cdot A_x^{[L]}$ have reductions
    $(V_{g,x},W_{g,x})$ onto $C_{g,x}=A_{gx}^{[L]}$ satisfying
    \begin{align}
        V_{g,x}B_{g,x}^{\mathbf w}W_{g,x}
         & =C_{g,x}^{\mathbf w},
        \label{eq:mpug_action_interior}            \\
        B_{g,x}^{\mathbf p}W_{g,x}C_{g,x}^{\mathbf c}
        V_{g,x}B_{g,x}^{\mathbf q}
         & =B_{g,x}^{\mathbf p\mathbf c\mathbf q}.
        \label{eq:mpug_action_exterior}
    \end{align}
    The interior word $\mathbf w$ may be empty. In the exterior identity
    $\mathbf c$ is nonempty and $|\mathbf p|,|\mathbf q|\geq1$.
    These are the interior and exterior equations of
    \cite[lines~1782--1873]{FrancoRubio2025MPUGauging}, with
    $A^\Gamma_{g,x}=V_{g,x}$ and
    $A^{\text{\reflectbox{$\Gamma$}}}_{g,x}=W_{g,x}$.
    The periodic action hypothesis is explicit, as in
    \cite[line~1064]{GarreRubio2023MPOSym}; see also
    \cite{gap:fbc25_mpu_action_stabilizer_source_audit}.

    In particular, for a transitive action of finite $G$ on $\mathsf X$
    and injective $A_x$, the label set is automatically finite. If
    $\{O_g\}$ is a simple injective MPU representation, this same blocking
    preserves that property, the injectivity of every $A_x$, and the
    periodic action law.
\end{lemma}

\begin{proof}\leanok
    \uses{def:mpug_action_buffer, lem:mpug_action_blocking,
        thm:mpug_permuted_block_l_symbol_spec, thm:mps_reduction_exterior_buffer}
    Choose the action reductions and the common length of
    Definition~\ref{def:mpug_action_buffer}. Their residual bounds follow
    from equality of the positive-length periodic vectors. Since
    $N_0(g,x)-1\leq L$, one exterior block supplies the required buffer
    for every pair. The interior identity holds on all words by the
    reduction equations. If $\mathsf X$ is nonempty and the action is
    transitive, the map $g\mapsto gx_0$ is onto, so $\mathsf X$ is finite;
    the empty case needs no choice of $x_0$. Positive blocking preserves
    injectivity and the MPU representation laws.
\end{proof}

\begin{lemma}[Reciprocal scalar action gauge]
    \label{lem:mpug_action_scalar_gauge}
    \lean{MPOTensor.GroupFamily.BlockActionData.smul,
        MPOTensor.GroupFamily.BlockActionData.hasExteriorBufferLength_smul,
        MPOTensor.GroupFamily.BlockActionData.block_smul}
    \leanok
    \uses{lem:mpug_action_common_blocking}
    For any nonzero scalars $\gamma_{g,x}$ the replacements
    $W_{g,x}\mapsto\gamma_{g,x}W_{g,x}$ and
    $V_{g,x}\mapsto\gamma_{g,x}^{-1}V_{g,x}$ preserve
    (\ref{eq:mpug_action_interior}) and
    (\ref{eq:mpug_action_exterior}) with the same exterior buffer length.
    They commute with physical blocking. This is the scalar action in
    \cite[lines~1871--1873]{FrancoRubio2025MPUGauging}.
\end{lemma}

\begin{proof}\leanok
    The two scalars cancel inside $V_{g,x}B_{g,x}^{\mathbf w}W_{g,x}$
    and $W_{g,x}C_{g,x}^{\mathbf c}V_{g,x}$. Blocking changes the physical
    tensors and leaves the reduction matrices unchanged.
\end{proof}
